\section{An exact local count identity and a two-code decomposition}
\label{sec:neighbor-counts}

The preceding pivot suggests counting over code neighbors, rather
than asking for an extension preserving every row. The following
identity makes that approach quantitative. Here $p$ is an odd prime,
$N=4p$, $m=N/2$, $C=C^\perp$, and $G\subsetneq C$ has dimension
$s=m-h$, $h\ge1$. Put
\[
 r_C=|S(C)|,\quad b_G=|S(G)|,\quad
 M_G=|\{v\in\{\pm1\}^N:v_1=1,\ v\bmod p\in G^\perp\}|.
\]
Let $\mathcal N_G(C)$ be the self-dual codes $C'\ne C$ with
$\dim(C\cap C')=m-1$ and $G\subseteq C'$.

\begin{theorem}[Exact restricted-neighbor incidence count]
\label{code:neighbor-average}
There are $(p^h-1)/(p-1)$ codes in $\mathcal N_G(C)$, and
\[
 \boxed{\quad
 \sum_{C'\in\mathcal N_G(C)}\bigl(|S(C')|-r_C\bigr)
 =M_G-r_C-(r_C-b_G)p^{h-1}.\quad}
\]
In particular, if the right-hand side is positive, some neighbor
strictly increases the sign count while preserving $G$.
\end{theorem}
\begin{proof}
A hyperplane $D$ of $C$ containing $G$ has $D^\perp/D$ a nondegenerate
two-dimensional hyperbolic space. It has exactly two isotropic lines,
one giving $C$ and the other a unique neighbor $C'$.
To see the last assertion directly, choose isotropic $e$ on the first
line and $f$ with $e\cdot f=1$, then replace $f$ by
$f-(f\cdot f)e/2$. In the resulting basis the norm is $2xy$,
so the isotropic lines are precisely the two coordinate axes.
Counting hyperplanes of $C/G$ gives the stated neighbor count.

Each normalized sign vector outside $C$ but in $G^\perp$ is isotropic,
since its norm is $N=0$ modulo $p$. It belongs to exactly one of
these neighbors, namely
$(C\cap v^\perp)+\langle v\rangle$. Thus the total number of
new sign classes, counted over the neighbors, is $M_G-r_C$.
An old sign class in $G$ is never lost. An old class outside $G$
is excluded by precisely $p^{h-1}$ hyperplanes containing $G$,
the difference between the hyperplane counts in dimensions
$h$ and $h-1$. Summing gains minus losses proves the identity.
\end{proof}

This is a finite incidence identity in the classical geometry of
self-dual code neighbors. It is not a general positivity theorem.
In particular, at $G=0$ its right-hand side is
$2^{N-1}-r_C(1+p^{m-1})$, which can be negative even when the
code has far fewer than $N$ sign classes. Unrestricted averaging
therefore does not guarantee a useful move.

\begin{corollary}[A modular pool splits into two codes]
\label{code:two-codes}
If $G$ has dimension $m-1$, there is exactly one other self-dual
code $C'$ containing $G$, and
\[
 \{ \text{sign classes in }G^\perp\}=S(C)\cup S(C'),\qquad
 S(C)\cap S(C')=S(G),\qquad
 |S(C')|=M_G-r_C+b_G.
\]
If $G$ contains five orthogonal sign rows, the sign words of each
code separately form a real orthogonal frame.
\end{corollary}
\begin{proof}
An isotropic residue in $G^\perp$ projects to one of the two
isotropic lines of $G^\perp/G$, hence belongs to $C$ or $C'$.
Their intersection is $G$. Count the union, and apply
Corollary~\ref{code:five} for real orthogonality.
\end{proof}

This removes a compatibility search in the codimension-one case.
For the terminal twelve-row order-twenty certificate, the ten rows
retained after deleting indices $3,8$ span such a $G$, and exact
enumeration gives
\[
 r_C=12,\qquad b_G=10,\qquad M_G=22,\qquad |S(C')|=20.
\]
The twelve exact extensions of the retained frame consist of the two
deleted old classes and the ten new classes. The two-code theorem,
rather than an independent guess that all candidates are compatible,
explains the resulting completion.

The checker also verifies the incidence identity for all four
neighbors containing a specified four-dimensional subspace of a
self-dual ternary code at order twelve. These are finite examples.
An all-order lower bound on the modular pool, or a guarantee that
some choice of retained space makes the displayed surplus positive,
has not been established.
